% Part: computability
% Chapter: machines-computations
% Section: configuration

\documentclass[../../../include/open-logic-section]{subfiles}

\begin{document}

\olfileid{cmp}{tur}{con}
\olsection{स्थितिवर्णने आणि संगणने}

\begin{explain}
मागे सम यंत्राच्या स्थितिवर्णनांचा क्रम कसा पाहिला, ते आठवा.
फीत, वाचन-लेखन शीर्ष आणि ट्यूरिंग यंत्राचा कार्यक्रम यांची बनलेली
काल्पनिक यंत्रणा ही ट्यूरिंग यंत्राचे संगणन कसे असते ते डोळ्यांसमोर
आणण्याचा केवळ अंतर्ज्ञानी मार्ग आहे. औपचारिकरीत्या दिलेल्या आदानावर
ट्यूरिंग यंत्राचे संगणन \emph{स्थितिवर्णनांची} क्रमिका म्हणून
परिभाषित करता येते. प्रत्येक स्थितिवर्णनात चिन्हांची एक क्रमिका
(संगणनाच्या त्या टप्प्यावर फितीतील मजकुराशी संबंधित), वाचन-लेखन
शीर्षाची जागा दर्शवणारी संख्या आणि एक अवस्था असते. त्यांच्या
आधारे ट्यूरिंग यंत्र~$M$ दिलेल्या आदानावर काय काढते हे परिभाषित
करता येते.
\end{explain}

\begin{defn}[स्थितिवर्णन]
ट्यूरिंग यंत्र $M = \tuple{Q, \Sigma, q_0,
\delta}$ चे \emph{स्थितिवर्णन} म्हणजे $\tuple{C, m, q}$ ही त्रयी;
येथे
\begin{enumerate}
\item $C \in \Sigma^*$ ही $\Sigma$ मधील चिन्हांची सांत क्रमिका आहे,
\item $m \in \Nat$ ही $< \len{C}$ असलेली संख्या आहे, आणि
\item $q \in Q$ आहे.
\end{enumerate}
अंतर्ज्ञानी अर्थाने क्रमिका~$C$ म्हणजे फितीतील मजकूर (सर्वांत
डावीकडील घरापासून शेवटच्या रिकाम्या नसलेल्या किंवा आधी भेट दिलेल्या
घरापर्यंतच्या सर्व घरांतील चिन्हे); $m$~म्हणजे वाचन-लेखन शीर्ष ज्या
घरावर आहे त्या घराचा क्रमांक (सर्वांत डावीकडील घराचा क्रमांक $0$
धरून); आणि $q$ म्हणजे यंत्राची सध्याची अवस्था.
\end{defn}

\begin{explain}
ट्यूरिंग यंत्राचे संभाव्य आदान ही चिन्हांची एक क्रमिका असते; बहुतेक
वेळा ती एखाद्या संख्येला काही रूपात सांकेतिक करणारी क्रमिका असते.
ट्यूरिंग यंत्राचे प्रारंभिक स्थितिवर्णन म्हणजे त्या आदानावर यंत्र
चालवायला सुरू करतानाचे स्थितिवर्णन: फितीत डाव्या टोकाचे चिन्ह असते
आणि त्याच्या लगेच उजवीकडील घरांत आदान लिहिलेले असते; वाचन-लेखन
शीर्ष आदानाच्या सर्वांत डावीकडील घरावर असते (म्हणजे डाव्या टोकाच्या
चिन्हाच्या उजवीकडील घरावर); आणि यंत्रणा निर्दिष्ट प्रारंभिक
अवस्थेत~$q_0$ असते.
%READERNOTE{OLTUR-002 — या औपचारिक प्रारंभिक स्थितिवर्णनात शीर्ष आदानाच्या पहिल्या घरावर आहे; आधीच्या परिचयात ते फितीच्या डाव्या टोकाच्या घरावर असल्याचे म्हटले आहे. दोन्ही गोठवलेली कथने जतन करून येथे सूत्राशी सुसंगत विधान दिले आहे.}
\end{explain}

\begin{defn}[प्रारंभिक स्थितिवर्णन]
$I \in \Sigma^*$ हे आदान असताना $M$ चे \emph{प्रारंभिक स्थितिवर्णन}
असे आहे:
\[
\tuple{\TMendtape \frown I, 1, q_0}.
\]
रिकामे आदान घेतल्यास फितीवरील पहिल्या रिकाम्या घराचाही
स्थितिवर्णनात समावेश करावा; अन्यथा शीर्षाचा क्रमांक आणि क्रमिकेची
लांबी यांच्यावरील वरची अट पूर्ण होत नाही.
%READERNOTE{OLTUR-004 — गोठवलेली व्याख्या रिकाम्या आदानालाही परवानगी देते; पण दाखवलेल्या त्रयीत शीर्षाचा क्रमांक एक आणि फितीच्या क्रमिकेची लांबीही एक असल्याने आधीची काटेकोर लहानपणाची अट मोडते. मराठीत रिकामे घर समाविष्ट करण्याची मर्यादित स्पष्टता दिली आहे; सूत्र जतन आहे.}
\end{defn}

\begin{explain}
$\frown$ हे चिन्ह \emph{जोडणीसाठी} आहे---आदान-चिन्हमाला डाव्या
टोकाच्या चिन्हाच्या लगेच उजवीकडे सुरू होते.
%READERNOTE{OLTUR-003 — गोठवलेल्या इंग्रजीत या स्पष्टीकरणात आदान डाव्या टोकाच्या चिन्हाच्या डावीकडे सुरू होते असे शब्दशः येते; लगेच आधीचे सूत्र आणि परिच्छेद उजवी बाजू निश्चित करतात. मराठीत उजवीकडे असे स्पष्ट केले आहे.}
\end{explain}

\begin{defn}
स्थितिवर्णन $\tuple{C, m, q}$ पासून \emph{एका पायरीत
$\tuple{C', m', q'}$ हे स्थितिवर्णन मिळते} ($M$ नुसार), असे आपण
तेव्हा आणि केवळ तेव्हाच म्हणतो, जेव्हा
\begin{enumerate}
\item $C$ चे $m$-वे चिन्ह $\sigma$ असते,
\item $M$ च्या सूचनासंचात $\delta(q, \sigma) =
  \tuple{q', \sigma', D}$ निर्दिष्ट असते,
\item $C'$ चे $m$-वे चिन्ह $\sigma'$ असते, आणि
\item
\begin{enumerate}
\item $D = L$ आणि $m>0$ असल्यास $m' = m - 1$ असते; अन्यथा $m'=0$ असते, किंवा
\item $D = R$ आणि $m' = m + 1$ असते, किंवा
\item $D = N$ आणि $m' = m$ असते,
\end{enumerate}
\item $m' = \len{C}$ असेल, तर $\len{C'} = \len{C} + 1$ असते आणि
  $C'$ चे $m'$-वे चिन्ह~$\TMblank$ असते. अन्यथा $\len{C'}=\len{C}$ असते.
\item $i < \len{C}$ आणि $i \neq m$ असलेल्या सर्व $i$ साठी $C'(i) = C(i)$ असते,
\end{enumerate}
\end{defn}

\begin{defn}\ollabel{defn:run-output}
\emph{आदान~$I$ वरील $M$ ची धाव} म्हणजे $M$ च्या स्थितिवर्णनांची
$C_i$ ही क्रमिका; येथे $C_0$ हे आदान~$I$ साठी $M$ चे प्रारंभिक
स्थितिवर्णन असते आणि प्रत्येक $C_i$ पासून एका पायरीत $C_{i+1}$
मिळते.

$C_k = \tuple{C, m, q}$ असेल, $C$ चे $m$-वे चिन्ह~$\sigma$ असेल
आणि $\delta(q,
\sigma)$ अपरिभाषित असेल, तर $M$ \emph{आदान $I$ वर $k$
पायऱ्यांनंतर थांबते} असे म्हणतो. त्या वेळी आदान~$I$ साठी $M$ चे
\emph{प्रदान}~$O$ असते; येथे $O$ ही शेवटी~$\TMblank$ नसलेली
चिन्हमाला असून एखाद्या~$j \in \Nat$ साठी
$C = \TMendtape \concat O \concat \TMblank^j$ असते.
($\TMblank^j$ म्हणजे $j$ $\TMblank$ चिन्हांची क्रमिका.)
ही प्रदानाची अट डाव्या टोकाचे चिन्ह शेवटपर्यंत शाबूत राहिलेल्या
धावांसाठीच थेट लागू होते. ते चिन्ह बदलून लिहिले असल्यास प्रदान
ठरवण्यासाठी अतिरिक्त नियम लागेल.
%READERNOTE{OLTUR-005 — आधीच्या विभागात डाव्या टोकाचे चिन्ह बदलून लिहिण्यास परवानगी आहे; गोठवलेल्या प्रदान-सूत्रात अंतिम फीत मात्र त्या चिन्हानेच सुरू होत असल्याचे गृहीत आहे. मराठीत या मर्यादेची स्पष्ट नोंद केली आहे.}
\end{defn}

\begin{explain}
या व्याख्येनुसार $M$ चे प्रदान~$O$ नेहमी~$\TMblank$ व्यतिरिक्त
इतर चिन्हाने संपते; किंवा $k$~व्या वेळी संपूर्ण फीत (सर्वांत
डावीकडील~$\TMendtape$ सोडून)~$\TMblank$ चिन्हांनी भरलेली असेल,
तर $O$~ही रिकामी चिन्हमाला असते.
\end{explain}

\end{document}
